\section{Matching the leading transports}\label{sec:transport-matching}

Proposition~\ref{prop:transport-matching} constructs the unique holomorphic
comparison of the two leading transports. We first derive a transport normal
form and prove a supported estimate for parameter dependence.
Throughout this section, a dot
between complex vectors denotes the complex bilinear Euclidean product.

\begin{lemma}[Transport normal form]\label{lem:transport-normal-form}
For one coefficient pair, set $g=\nabla\log m$. The change of variables
\begin{equation}\label{eq:transport-change}
 p=\sqrt m\,a,
 \qquad
 s=-\frac{k}{2\sqrt m}\,b-\frac12 g\cdot p
\end{equation}
is a smooth invertible map of $\Ecal_\theta$ to itself. It transforms
\eqref{eq:constrained-transport} into
\begin{equation}\label{eq:transport-matrix}
 D_\theta\binom{p}{s}
   =M_\theta\binom{p}{s},
 \qquad
 M_\theta=
 \begin{pmatrix}
  0&\theta\\
  Q_\theta\cdot&T_\theta
 \end{pmatrix}
 \quad\text{on }H_\theta\oplus\C,
\end{equation}
where
\begin{equation}\label{eq:transport-coefficients}
 \begin{aligned}
  T_\theta&=D_\theta\log(k/\ell),\\
  Q_\theta&=\frac12\left[
   \left(\frac m\ell D_\theta\log k
          -\frac12\theta\cdot g\right)g
         -\frac m\ell D_\theta g\right].
 \end{aligned}
\end{equation}
In particular, the bottom-left entry in \eqref{eq:transport-matrix}
is the covector $p\mapsto Q_\theta\cdot p$. Both $D_\theta$ and
$M_\theta$ depend linearly and holomorphically on $\theta$.
\end{lemma}

\begin{proof}
The inverse of \eqref{eq:transport-change} is
\[
 a=m^{-1/2}p,
 \qquad
 b=-\frac{2\sqrt m}{k}\left(s+\frac12g\cdot p\right).
\]
It is well defined because $m,k>0$, and multiplication by $\sqrt m$
preserves $H_\theta$. Write $D=D_\theta$. Substitution into the first
equation in \eqref{eq:leading-operators} gives
\[
 0=2\sqrt m\,Dp+\theta(kb+\sqrt m\,g\cdot p),
 \qquad\text{hence}\qquad Dp=\theta s.
\]
Dividing the second equation by two and using
$\nabla m=mg$ yields
\[
 0=\ell Db+(Dk)b+2\nabla m\cdot Da+(D\nabla m)\cdot a.
\]
Here
\[
 \begin{aligned}
 2\nabla m\cdot Da
   &=2\sqrt m\,g\cdot Dp
     -\sqrt m\,(D\log m)g\cdot p,\\
 (D\nabla m)\cdot a
   &=\sqrt m\,\bigl((D\log m)g+Dg\bigr)\cdot p.
 \end{aligned}
\]
The terms containing $(D\log m)g\cdot p$ cancel. Thus
\begin{equation}\label{eq:transformed-divergence-transport}
 \ell Db=-(Dk)b
   -\sqrt m\bigl(2(\theta\cdot g)s+(Dg)\cdot p\bigr).
\end{equation}
Differentiate the definition of $s$ and substitute
\eqref{eq:transformed-divergence-transport}, the expression for $b$,
and $Dp=\theta s$. With $c=m/\ell$ and $\gamma=D\log m=\theta\cdot g$,
the result is
\[
 Ds=c(D\log k-\gamma)s
   +\frac12\bigl[(cD\log k-\gamma/2)g-cDg\bigr]\cdot p.
\]
Finally, since $\ell=k+m$,
\[
 D\log(k/\ell)
 =D\log k-\frac{Dk+Dm}{\ell}
 =\frac m\ell(D\log k-D\log m).
\]
This proves the formulas. The top row takes values in $H_\theta$
because $\theta\cdot\theta=0$. All the direction dependence displayed
in the formulas is linear.
\end{proof}

The comparison equation can now be described in terms of the columns
that the physical construction must supply. Fix $\theta$ and a cylinder
$D_{2r}^2\times J$. Suppose for the moment that
$Y_1:\C^3\to\Ecal_\theta$ is a smooth invertible matrix of normalized
transport columns for the first medium, and that
$Y_2:\C^3\to\Ecal_\theta$ has locally $L^2$ columns satisfying the
second transport equation in distributions:
\[
 D_\theta Y_j=M_{j,\theta}Y_j,\qquad j=1,2.
\]
The columns in each $Y_j$ are normalized by that medium's own change
of variables \eqref{eq:transport-change}. These changes agree outside
$\overline\Omega$, where the extended coefficients agree. Thus exterior
agreement of the corresponding physical leading columns gives
$Y_2=Y_1$ there.

Under this exterior agreement, define
\[
 G=Y_2Y_1^{-1}\in\operatorname{End}(\Ecal_\theta).
\]
Only $Y_1$ is required to be invertible. Multiplication by its smooth
inverse is legitimate for locally $L^2$ columns, and the product rule
in distributions gives
\[
 D_\theta G=M_{2,\theta}G-GM_{1,\theta},
 \qquad G=\Id\quad\text{outside }\overline\Omega
\]
on the cylinder. Consequently $W=G-\Id$ is locally $L^2$ and supported
in $\overline\Omega$ within that cylinder.

For a transverse coordinate $t$, restrict the coefficients to the
corresponding plane and use a basis of $\Ecal_\theta$ fixed in $x$.
On planar endomorphism-valued unknowns define
\begin{equation}\label{eq:planar-comparison-operator}
 \Pcal_{\theta,t}W
      =D_\theta W-M_{2,\theta}W+WM_{1,\theta},
 \qquad
 f_{\theta,t}=M_{2,\theta}-M_{1,\theta}.
\end{equation}
The supported planar equation required for the comparison is
\begin{equation}\label{eq:inhomogeneous-comparison}
 \Pcal_{\theta,t}W=f_{\theta,t}.
\end{equation}
Matrices here may be regarded as vectors of their nine entries.
For almost every $t$, the $L^2$ planar restriction of the local $W$
is supported in $\overline{D_r^2}$ because $\overline\Omega\subset B_r$.
The physical transfer and weak-limit argument below will supply the
second columns $Y_2$. Slicing their distributional cylinder equation
will initially give \eqref{eq:inhomogeneous-comparison} only for almost
every $t$. We therefore need both uniqueness for supported planar
solutions and a way to pass from this almost-everywhere construction
to smooth dependence on every parameter. The next lemma provides these
facts on fixed Hilbert spaces, which will also allow the planes to vary.

\begin{lemma}[Supported planar transport]\label{lem:supported-transport}
Let $\mathcal U\subset\R^d$ be open, let $n\geq1$, and suppose that
\[
 \Pcal_\rho
   =\gamma(\rho)(\partial_{y_1}+i\partial_{y_2})
       +A(y,\rho),
 \qquad \rho\in\mathcal U,
\]
where $\gamma$ is smooth and nonzero and $A$ is a smooth complex
$n\times n$ matrix on a neighborhood of
$\overline{D_{2r}^2}\times\mathcal U$. Then the following statements hold.
\begin{enumerate}[label=\textup{(\roman*)}]
\item A distributional solution of $\Pcal_\rho W=0$ on $D_{2r}^2$
whose support is a compact subset of that disk is zero.
\item For every integer $j\geq0$ and every
$W\in H^{j+1}(D_{2r}^2;\C^n)$ supported in $\overline{D_r^2}$,
\begin{equation}\label{eq:supported-transport-estimate}
 \norm{W}_{H^{j+1}(D_{2r}^2)}
 \leq C_j\norm{\Pcal_\rho W}_{H^j(D_{2r}^2)},
 \qquad
 \supp W\subset\overline{D_r^2}.
\end{equation}
The constants are locally uniform in $\rho$.
\item Suppose that $f(y,\rho)$ is smooth on a neighborhood of
$\overline{D_{2r}^2}\times\mathcal U$. If the equation
$\Pcal_\rho W=f(\cdot,\rho)$ has an $L^2$ solution supported in
$\overline{D_r^2}$ for a dense set of parameters $\rho$, then it has
such a solution for every $\rho$. These solutions are unique and form
a smooth function of $(y,\rho)$.
\end{enumerate}
The same conclusions apply to matrix-valued unknowns, after their
entries are regarded as a vector.
\end{lemma}

\begin{proof}
For a fixed parameter, $\Pcal_\rho$ is elliptic in the two real planar
variables: its principal symbol is
$\gamma(\rho)(i\xi_1-\xi_2)$, which is nonzero for
$\xi\in\R^2\setminus\{0\}$. On each sufficiently small disk,
Lemma~\ref{lem:planar-frames} gives a smooth invertible matrix $F$
satisfying $\Pcal_\rho F=0$. If $\Pcal_\rho W=0$, then
\[
 (\partial_{y_1}+i\partial_{y_2})(F^{-1}W)=0
\]
in distributions. Its components are therefore holomorphic functions
of $y_1+iy_2$. A homogeneous solution that vanishes on an open set
vanishes on every overlapping frame disk by the holomorphic identity
theorem. Connectedness propagates this conclusion throughout
$D_{2r}^2$. A compactly supported solution vanishes near the boundary
of the disk, proving \textup{(i)}.

We give the estimate on the fixed supported space. Extending $W$ by
zero to $\R^2$ is harmless because its support lies in $\overline{D_r^2}$.
The Fourier symbol just computed gives
\[
 \norm{W}_{H^{j+1}}
 \leq C\bigl(
    \norm{\gamma(\rho)(\partial_{y_1}+i\partial_{y_2})W}_{H^j}
       +\norm{W}_{L^2}\bigr).
\]
Here and below in this proof the norms can equivalently be taken on
$D_{2r}^2$, since the functions under consideration are supported in
the smaller disk. Smooth multiplication and, for $j\geq1$, the
interpolation estimate
$\norm{W}_{H^j}\leq\varepsilon\norm{W}_{H^{j+1}}
 +C_\varepsilon\norm{W}_{L^2}$ give
\begin{equation}\label{eq:supported-elliptic-preestimate}
 \norm{W}_{H^{j+1}}
 \leq C\bigl(\norm{\Pcal_\rho W}_{H^j}
                  +\norm{W}_{L^2}\bigr).
\end{equation}
The coefficients can be cut off outside a neighborhood of the smaller
disk for this calculation. We will also use local regularity for an
$L^2$ distributional solution with smooth right side. In a local frame
$F$, its equation becomes a scalar Cauchy--Riemann equation componentwise
with smooth right side. Subtracting a smooth local particular solution,
as supplied by Lemma~\ref{lem:planar-frames}, leaves holomorphic
components. The original solution is therefore smooth in the disk.

To remove the last term in \eqref{eq:supported-elliptic-preestimate},
suppose that \eqref{eq:supported-transport-estimate} fails at a fixed
parameter. There would then be a sequence $W_\nu$, supported in
$\overline{D_r^2}$, such that
\[
 \norm{W_\nu}_{H^{j+1}}=1,
 \qquad
 \norm{\Pcal_\rho W_\nu}_{H^j}\longrightarrow0.
\]
After passing to a subsequence, compactness of the embedding into
$L^2(D_{2r}^2)$ gives a strong $L^2$ limit $W$. This limit is supported
in $\overline{D_r^2}$ and solves $\Pcal_\rho W=0$ distributionally.
Part~\textup{(i)} gives $W=0$. Equation
\eqref{eq:supported-elliptic-preestimate} now contradicts the
normalization of $W_\nu$. This proves \textup{(ii)} at a fixed parameter.

For the parameter assertions, equip the fixed complex Hilbert spaces
\[
 \begin{aligned}
 X_j&=\{W\in H^{j+1}(D_{2r}^2;\C^n):
                  \supp W\subset\overline{D_r^2}\},\\
 Z_j&=H^j(D_{2r}^2;\C^n)
 \end{aligned}
\]
with their usual Hermitian Sobolev inner products. The subspace $X_j$
is closed. The map $\rho\mapsto\Pcal_\rho$ is smooth in
$\mathcal L(X_j,Z_j)$. The estimate at $\rho_0$ persists in a
neighborhood: absorb
$\norm{(\Pcal_\rho-\Pcal_{\rho_0})W}_{Z_j}$ into its left side when
the operator norm of the difference is small. This proves local
uniformity in \textup{(ii)}.

Write $\Pcal_\rho^*:Z_j\to X_j$ for the Hilbert space adjoint with
respect to these fixed inner products. In particular, this is not an
adjoint in the complex bilinear pairing used for the null vectors.
The lower bound in \textup{(ii)} implies that
$\Pcal_\rho^*\Pcal_\rho:X_j\to X_j$ is coercive and invertible.
Consequently
\begin{equation}\label{eq:supported-left-inverse}
 T_{\rho,j}
   =(\Pcal_\rho^*\Pcal_\rho)^{-1}\Pcal_\rho^*
       :Z_j\longrightarrow X_j
\end{equation}
is a bounded left inverse. It depends smoothly on $\rho$, since
adjoints and inversion of bounded invertible operators do so locally.
This argument uses a lower bound for $\Pcal_\rho$; it requires no
surjectivity onto $Z_j$.

For each $j$, set $W_j(\rho)=T_{\rho,j}f(\cdot,\rho)$. An $L^2$
supported solution, whenever it exists, is smooth by interior
ellipticity and hence belongs to every $X_j$. It must equal
$W_j(\rho)$ by the left-inverse identity. Thus the smooth
$Z_j$-valued residual
\[
 \Pcal_\rho W_j(\rho)-f(\cdot,\rho)
\]
vanishes on the assumed dense set, and continuity makes it zero for
every $\rho$. Uniqueness in \textup{(i)} identifies the solutions for
different $j$. Smooth dependence with values in every $H^{j+1}$,
together with Sobolev embedding in the planar variables, gives joint
smoothness in $(y,\rho)$. This proves \textup{(iii)}.
\end{proof}

Let
\[
 \mathcal C=\{[\theta]\in\C\mathrm P^2:
                         \theta\cdot\theta=0\}
\]
be the projective null conic. The fibers
$\Ecal_\theta=H_\theta\oplus\C$ form a holomorphic subbundle
$\Ecal$ of the trivial bundle $\mathcal C\times\C^4$. For example,
on a chart where $\theta_i\ne0$, the vectors
$e_j-(\theta_j/\theta_i)e_i$, $j\ne i$, give a holomorphic frame
of $H_\theta$; adjoining the scalar component gives one of $\Ecal$.

Ceki\'c constructs a transport comparison for connection Laplacians
that equals the identity in the outer exterior component
\cite[Theorem~3.4 in the arXiv version]{Cekic2025}, using full
connection-system boundary data on a trivial vector bundle. The comparison
below is obtained from physical elastic transfer and weak limits on the
constrained fibers $\Ecal_\theta$.

\begin{proposition}[Matching of the transports]\label{prop:transport-matching}
For the two extended coefficient pairs of
Lemma~\ref{lem:common-extension}, equality of the physical
displacement-to-traction maps implies that there is a unique smooth
field
\[
 G(x,[\theta])\in\operatorname{End}(\Ecal_\theta),
 \qquad (x,[\theta])\in\R^3\times\mathcal C,
\]
satisfying, for every nonzero representative $\theta$,
\begin{equation}\label{eq:comparison-transport}
 \begin{aligned}
 D_\theta G&=M_{2,\theta}G-GM_{1,\theta},\\
 G(x,[\theta])&=\Id\qquad\text{if }|x|>2r.
 \end{aligned}
\end{equation}
For each fixed $x$, this field is a holomorphic section of
$\operatorname{End}(\Ecal)$ over $\mathcal C$. In fact, for each
fixed direction the first equation has at most one smooth solution whose
difference from the identity has compact spatial support.
\end{proposition}

\begin{proof}
\emph{Physical transfer and weak limits.}
Fix a nonzero null vector $\theta$ and $t_0\in(-2r,2r)$. Apply
Proposition~\ref{prop:physical-amplitudes} to the first coefficient
pair. It supplies three exact smooth solutions on $B_{5r}$, with the
limits and estimates in \eqref{eq:physical-leading-limit}, and an
interval $J$ containing $t_0$ on which their leading columns form a
frame on $D_{2r}^2\times J$. We perform the following construction
for each column, temporarily suppressing its index.

By Proposition~\ref{prop:physical-transfer}, the first solution
$u_{1,\tau}$ has a smooth transferred solution $u_{2,\tau}$ on
$B_{5r}$, with
$u_{2,\tau}=u_{1,\tau}$ outside $\overline\Omega$. Define the
physical augmented vectors
\[
 V_{j,\tau}=(u_{j,\tau},\diverg u_{j,\tau}),\qquad j=1,2.
\]
The difference $V_{2,\tau}-V_{1,\tau}$ is smooth and compactly
supported in $\overline\Omega\subset B_r$. By
Lemma~\ref{lem:augmented-system},
\begin{equation}\label{eq:physical-transfer-forcing}
 \Lcal_2(V_{2,\tau}-V_{1,\tau})
       =(\Lcal_1-\Lcal_2)V_{1,\tau}.
\end{equation}
Both normalized systems have principal part $\Delta\Id_4$.
Their difference is therefore an operator of order at most one,
whose coefficients are supported in $\overline\Omega$.
The weighted derivative bound in
\eqref{eq:physical-leading-limit} implies
\[
 \begin{aligned}
 \norm{E_\tau^{-1}(\Lcal_1-\Lcal_2)V_{1,\tau}}_{L^2}
 &\leq C\left(
     \norm{E_\tau^{-1}V_{1,\tau}}_{L^2(B_{4r})}
    +\norm{E_\tau^{-1}\nabla V_{1,\tau}}_{L^2(B_{4r})}
          \right)\\
 &=O(\tau).
 \end{aligned}
\]
The function on the left is extended by zero off its compact support
when a norm on $\R^3$ is used. Apply
Proposition~\ref{prop:carleman} to the difference in
\eqref{eq:physical-transfer-forcing}. Since
$\norm{F}_{H_h^{-1}}\leq\norm{F}_{L^2}$ and $h=\tau^{-1}$, we obtain
\begin{equation}\label{eq:bounded-transferred-amplitudes}
 \begin{aligned}
 \norm{E_\tau^{-1}(V_{2,\tau}-V_{1,\tau})}_{H_h^1}
 &\leq Ch\norm{E_\tau^{-1}(\Lcal_1-\Lcal_2)V_{1,\tau}}_{H_h^{-1}}\\
 &\leq C.
 \end{aligned}
\end{equation}
Thus $E_\tau^{-1}V_{2,\tau}$ is bounded in $L^2(B_{4r})$.
After passage to a subsequence, chosen simultaneously for the three
columns, it converges weakly to a vector $(\widetilde a,\widetilde b)$.
The strong convergence of the first amplitudes and the exterior
agreement of the physical solutions give
\begin{equation}\label{eq:exterior-leading-agreement}
 (\widetilde a,\widetilde b)=(a_0,b_0)
 \quad\text{almost everywhere on }
 B_{4r}\setminus\overline\Omega.
\end{equation}

These weak limits satisfy the constrained transport for the second
pair. To see this directly, write
$(a_\tau,b_\tau)=E_\tau^{-1}V_{2,\tau}$. The exact equations and
the null condition give
\[
 \begin{aligned}
 R_2^\theta(a_\tau,b_\tau)
       +\tau^{-1}R_2(a_\tau,b_\tau)&=0,\\
 S_2^\theta(a_\tau,b_\tau)
       +\tau^{-1}S_2(a_\tau,b_\tau)&=0.
 \end{aligned}
\]
When tested against a fixed compactly supported smooth function,
the terms multiplied by $\tau^{-1}$ tend to zero: all their
derivatives can be transferred to that test function, while the
amplitudes remain bounded in $L^2$. Weak convergence passes the
remaining terms to the limit. Moreover, the physical divergence
identity says
\[
 \theta\cdot a_\tau
   =\tau^{-1}(b_\tau-\diverg a_\tau),
\]
so $\theta\cdot\widetilde a=0$ distributionally. We have proved
\[
 R_2^\theta(\widetilde a,\widetilde b)=0,
 \qquad S_2^\theta(\widetilde a,\widetilde b)=0,
 \qquad (\widetilde a,\widetilde b)\in\Ecal_\theta.
\]
No derivative estimate on the transferred amplitudes, beyond their
$L^2$ bound, is needed for this passage to the limit.

The physical transfer has therefore produced the second transport
columns, but only as weak limits. We next obtain a supported comparison
on almost every transverse plane and use supported uniqueness to extend
it to every plane.

\emph{A comparison on almost every plane.}
Apply \eqref{eq:transport-change} to the three leading columns for
each pair and denote the resulting maps by
$Y_j:\C^3\to\Ecal_\theta$. For $j=2$, these columns are initially
locally $L^2$; they satisfy the equations distributionally. Since
$J\subset(-2r,2r)$, the cylinder on which the first columns form a frame
satisfies
\[
 D_{2r}^2\times J\subset B_{\sqrt8\,r}\subset B_{4r}.
\]
Thus the weak limits constructed above are defined throughout this
cylinder, where
\[
 D_\theta Y_j=M_{j,\theta}Y_j,
\]
and $Y_1$ is smooth and invertible. The change of variables agrees
for the two pairs outside $\overline\Omega$. Define
\[
 G=Y_2Y_1^{-1}.
\]
The calculation preceding Lemma~\ref{lem:supported-transport} now applies:
$G$ is a locally $L^2$ endomorphism on the cylinder, equals the identity
outside $\overline\Omega$, and satisfies the first equation of
\eqref{eq:comparison-transport} distributionally. With $W=G-\Id$, this
is the cylinder version of \eqref{eq:inhomogeneous-comparison}.
The operator $D_\theta$ differentiates only the planar variables.
Testing the cylinder equation with products of a planar test
function and a transverse test function, and then using a countable
dense collection of planar tests, gives
\eqref{eq:inhomogeneous-comparison} on $D_{2r}^2$ for almost every
$t\in J$. For these $t$, $W(\cdot,t)$ belongs to $L^2(D_{2r}^2)$
and has support in $\overline{D_r^2}$. Indeed, $W=0$ when $|y|>r$,
and on its remaining support multiplication by $Y_1^{-1}$ is locally
bounded uniformly in $t$. These assertions can first be made on
compact subintervals of $J$ and then exhausted over $J$.

\emph{Existence on every plane.}
In the orthonormal coordinates associated with the fixed $\theta$,
\[
 D_\theta=|\operatorname{Re}\theta|
                 (\partial_{y_1}+i\partial_{y_2}).
\]
Thus \eqref{eq:planar-comparison-operator}, as an operator on matrix
entries, has the form of Lemma~\ref{lem:supported-transport}.
Its coefficients and the source $f_{\theta,t}$ are smooth in $(y,t)$.
The set of parameters for which a supported solution has just been
obtained has full measure in $J$, and hence is dense. Part~\textup{(iii)}
of that lemma extends existence to every $t\in J$ and makes $W$
smooth in $(y,t)$. It also makes the choice unique, so the results
from different intervals $J$ agree on their overlaps.

Since $t_0$ was arbitrary, this constructs $G$ for every
$-2r<t<2r$. Extend $W=G-\Id$ by zero in the planar variables
outside $D_{2r}^2$. This extension is smooth because $W$ is already
zero for $r<|y|<2r$. The equation holds across this extension, and
it holds farther out because $f_{\theta,t}=0$ there. If $|t|>r$,
the entire plane misses $\overline\Omega$, so $f_{\theta,t}=0$.
Part~\textup{(i)} of Lemma~\ref{lem:supported-transport} then gives
$W=0$ on that plane. Extending by zero also for $|t|\geq2r$ is
therefore smooth. For this fixed direction we have constructed a
global smooth comparison with
\begin{equation}\label{eq:comparison-uniform-support}
 \supp(G(\cdot,\theta)-\Id)
 \subset\{(y,t):|y|\leq r,\ |t|\leq r\}
 \subset\overline{B_{2r}}.
\end{equation}
The final containing ball is independent of the direction.

For later use, uniqueness holds even if a larger compact support is
allowed. The difference of two such comparisons solves the
homogeneous equation on every plane and has compact support there.
The local-frame and holomorphic-continuation argument in
Lemma~\ref{lem:supported-transport}, applied on a disk containing
that support, makes the difference zero on each plane.

\emph{Smooth dependence on the direction.}
The preceding construction gives a unique comparison separately for
every nonzero null vector $\theta$. If $c\in\C\setminus\{0\}$,
then $\Ecal_{c\theta}=\Ecal_\theta$, and
\[
 D_{c\theta}=cD_\theta,
 \qquad M_{j,c\theta}=cM_{j,\theta}.
\]
The equations for $\theta$ and $c\theta$ are consequently identical
after division by $c$. Compact-support uniqueness gives
$G(x,c\theta)=G(x,\theta)$, so the field is already well defined
on projective directions as a pointwise family.

Choose a local smooth representative $\theta(v)$ of those directions,
where $v$ ranges over an open subset of $\R^2$, and a smooth local
frame of $\Ecal$. Define the orthonormal vectors
\[
 e_1(v)=\frac{\operatorname{Re}\theta(v)}
                   {|\operatorname{Re}\theta(v)|},
 \quad
 e_2(v)=\frac{\operatorname{Im}\theta(v)}
                   {|\operatorname{Im}\theta(v)|},
 \quad
 e_3(v)=e_1(v)\times e_2(v).
\]
They vary smoothly: the real and imaginary parts of a nonzero null
vector are nonzero, perpendicular, and of equal length. In the
coordinates
\[
 x=y_1e_1(v)+y_2e_2(v)+te_3(v),
\]
and the chosen bundle frame, \eqref{eq:inhomogeneous-comparison}
is a smooth family of operators of the form in
Lemma~\ref{lem:supported-transport}, with parameters $(t,v)$ on
the fixed disk $D_{2r}^2$. For every value of these parameters a
solution supported in $\overline{D_r^2}$ has already been
constructed. The smooth left inverse
\eqref{eq:supported-left-inverse} therefore expresses this unique
solution as a smooth function of $(y,t,v)$. Returning to physical
coordinates gives joint smoothness in $(x,v)$.

This reasoning uses the separate fixed-direction existence results
and their uniqueness. In particular, none of the earlier weakly
convergent subsequences has to be chosen simultaneously for varying
directions.

\emph{Holomorphic dependence at fixed physical coordinates.}
A close precedent for this parameter argument is Ceki\'c
\cite[Theorem~3.4 and Lemma~4.1 in the arXiv version]{Cekic2025}, whose
parameter argument follows Eskin's complex-parameter method \cite{Eskin2001}.
Now let $v$ be a complex local coordinate on $\mathcal C$. Choose
a holomorphic nonzero representative $\theta(v)$ and a holomorphic
frame of $\Ecal$ on this chart. This frame depends only on $v$.
In this frame the matrices $M_{j,\theta(v)}(x)$ are holomorphic in
$v$ for each fixed $x$: the ambient formulas
\eqref{eq:transport-matrix}--\eqref{eq:transport-coefficients} are
holomorphic in $\theta$, and passage to a holomorphic bundle frame
preserves that property.

Represent $G$ in the same frame and differentiate the first equation
of \eqref{eq:comparison-transport} by $\partial_{\bar v}$,
\emph{holding $x\in\R^3$ fixed}. Joint smoothness just established
justifies this differentiation. Since the representative and the
frame are holomorphic and the frame is independent of $x$, the
result is exactly
\[
 D_{\theta(v)}(\partial_{\bar v}G)
  -M_{2,\theta(v)}\partial_{\bar v}G
  +(\partial_{\bar v}G)M_{1,\theta(v)}=0.
\]
The derivative here is a matrix in a holomorphic trivialization,
hence represents an endomorphism of the same fiber $\Ecal_{\theta(v)}$.
The common support bound \eqref{eq:comparison-uniform-support}
implies that this derivative is zero for $|x|>2r$. Compact-support
uniqueness on each plane gives $\partial_{\bar v}G=0$. The smooth
moving coordinates were used to establish joint regularity; the
equation was differentiated only after returning to fixed physical
coordinates. We conclude that $G(x,\cdot)$ is holomorphic on
$\mathcal C$ for every $x$, completing the proof.
\end{proof}
